\nivelbasico

\begin{ejercicio}
Dos automóviles, \(A\) y \(B\), se desplazan sobre una carretera
recta. El automóvil \(A\) tiene una velocidad de
\(\SI{24.0}{\meter\per\second}\) hacia el este, mientras que
\(B\) se desplaza a \(\SI{16.0}{\meter\per\second}\) en
la misma dirección.

Determine:

\begin{enumerate}[label=\alph*)]
    \item La velocidad de \(A\) respecto de \(B\).
    \item La velocidad de \(B\) respecto de \(A\).
    \item La rapidez relativa entre ambos.
\end{enumerate}
\end{ejercicio}

\begin{respuesta}
\[
\boxed{
\vect{v}_{A/B}
=
+8.0\ihat\,
\si{\meter\per\second}
}
\]

\[
\boxed{
\vect{v}_{B/A}
=
-8.0\ihat\,
\si{\meter\per\second}
}
\]

\[
\boxed{
v_{\mathrm{rel}}
=
\SI{8.0}{\meter\per\second}
}
\]
\end{respuesta}

\begin{explicacion}
Elegimos el eje \(+x\) hacia el este.

Las velocidades respecto del suelo son:

\[
\vect{v}_A
=
24.0\ihat\,
\si{\meter\per\second},
\]

\[
\vect{v}_B
=
16.0\ihat\,
\si{\meter\per\second}.
\]

Por definición, la velocidad relativa de \(A\) respecto
de \(B\) es:

\[
\vect{v}_{A/B}
=
\vect{v}_A-\vect{v}_B.
\]

Sustituyendo:

\[
\vect{v}_{A/B}
=
(24.0-16.0)\ihat\,
\si{\meter\per\second}.
\]

Por tanto:

\[
\boxed{
\vect{v}_{A/B}
=
+8.0\ihat\,
\si{\meter\per\second}
}.
\]

Desde el automóvil \(B\), el automóvil \(A\)
parece avanzar hacia el este.

Ahora invertimos la referencia:

\[
\vect{v}_{B/A}
=
\vect{v}_B-\vect{v}_A.
\]

Entonces:

\[
\boxed{
\vect{v}_{B/A}
=
-8.0\ihat\,
\si{\meter\per\second}
}.
\]

El signo negativo indica movimiento relativo hacia el oeste.

La rapidez relativa es la magnitud del vector:

\[
v_{\mathrm{rel}}
=
\norm{\vect{v}_{A/B}}.
\]

Por tanto:

\[
\boxed{
v_{\mathrm{rel}}
=
\SI{8.0}{\meter\per\second}
}.
\]

Las velocidades relativas tienen igual magnitud
y sentidos opuestos.
\end{explicacion}


\begin{ejercicio}
Dos vehículos se aproximan sobre una carretera recta.

El vehículo \(A\) se desplaza hacia el este con rapidez
\(\SI{20.0}{\meter\per\second}\), mientras que \(B\)
se desplaza hacia el oeste con rapidez
\(\SI{15.0}{\meter\per\second}\).

Determine la velocidad de \(A\) respecto de \(B\)
y explique por qué su magnitud es mayor que
la rapidez de cualquiera de los vehículos.
\end{ejercicio}

\begin{respuesta}
\[
\boxed{
\vect{v}_{A/B}
=
+35.0\ihat\,
\si{\meter\per\second}
}
\]

Su magnitud es:

\[
\boxed{
v_{A/B}
=
\SI{35.0}{\meter\per\second}
}.
\]
\end{respuesta}

\begin{explicacion}
Tomamos \(+x\) hacia el este.

Por tanto:

\[
v_{Ax}=+\SI{20.0}{\meter\per\second},
\]

\[
v_{Bx}=-\SI{15.0}{\meter\per\second}.
\]

El signo negativo de \(v_{Bx}\) indica que \(B\)
se desplaza hacia el oeste.

La velocidad relativa se obtiene mediante:

\[
v_{A/B,x}=v_{Ax}-v_{Bx}.
\]

Entonces:

\[
v_{A/B,x}=20.0-(-15.0).
\]

Por tanto:

\[
\boxed{
v_{A/B,x}
=
+\SI{35.0}{\meter\per\second}
}.
\]

La magnitud es:

\[
\boxed{
v_{A/B}
=
\SI{35.0}{\meter\per\second}
}.
\]

Las rapideces aparecen sumadas porque las componentes
de velocidad tienen signos opuestos.

No hemos utilizado una regla especial para vehículos
que se aproximan, sino la misma relación vectorial
general.

La velocidad relativa indica el movimiento que
\(A\) presenta para un observador asociado a \(B\).
\end{explicacion}


\begin{ejercicio}
Una persona camina hacia la parte delantera de un tren
con rapidez \(\SI{1.50}{\meter\per\second}\) respecto
del tren.

El tren se desplaza hacia el este con velocidad
\(\SI{18.0}{\meter\per\second}\) respecto del suelo.

Determine la velocidad de la persona respecto del suelo
cuando:

\begin{enumerate}[label=\alph*)]
    \item Camina hacia la parte delantera.
    \item Camina hacia la parte trasera con la misma rapidez.
    \item Permanece en reposo respecto del tren.
\end{enumerate}
\end{ejercicio}

\begin{respuesta}
\[
\boxed{
\vect{v}_{P/S}
=
+19.5\ihat\,
\si{\meter\per\second}
}
\qquad \text{a)}
\]

\[
\boxed{
\vect{v}_{P/S}
=
+16.5\ihat\,
\si{\meter\per\second}
}
\qquad \text{b)}
\]

\[
\boxed{
\vect{v}_{P/S}
=
+18.0\ihat\,
\si{\meter\per\second}
}
\qquad \text{c)}
\]
\end{respuesta}

\begin{explicacion}
Elegimos \(+x\) hacia el este, que coincide con
el sentido de avance del tren.

La composición clásica de velocidades establece:

\[
\vect{v}_{P/S}
=
\vect{v}_{P/T}
+
\vect{v}_{T/S}.
\]

La velocidad del tren es:

\[
\vect{v}_{T/S}
=
18.0\ihat\,
\si{\meter\per\second}.
\]

\textbf{a) Camina hacia adelante.}

La velocidad de la persona respecto del tren es:

\[
\vect{v}_{P/T}
=
1.50\ihat\,
\si{\meter\per\second}.
\]

Entonces:

\[
\vect{v}_{P/S}
=
(1.50+18.0)\ihat\,
\si{\meter\per\second}.
\]

Por tanto:

\[
\boxed{
\vect{v}_{P/S}
=
19.5\ihat\,
\si{\meter\per\second}
}.
\]

\textbf{b) Camina hacia atrás.}

Ahora:

\[
\vect{v}_{P/T}
=
-1.50\ihat\,
\si{\meter\per\second}.
\]

Entonces:

\[
\vect{v}_{P/S}
=
(-1.50+18.0)\ihat\,
\si{\meter\per\second}.
\]

Así:

\[
\boxed{
\vect{v}_{P/S}
=
16.5\ihat\,
\si{\meter\per\second}
}.
\]

\textbf{c) Permanece en reposo dentro del tren.}

En este caso:

\[
\vect{v}_{P/T}=\vect{0}.
\]

Por tanto:

\[
\boxed{
\vect{v}_{P/S}
=
18.0\ihat\,
\si{\meter\per\second}
}.
\]

El resultado muestra que estar en reposo respecto
del tren no significa estar en reposo respecto del suelo.
\end{explicacion}


\begin{ejercicio}
En un sistema cartesiano fijo al suelo, dos partículas
se encuentran en las posiciones:

\[
\vect{r}_A
=
(12.0\ihat+8.0\jhat)\,\si{\meter},
\]

\[
\vect{r}_B
=
(4.0\ihat+14.0\jhat)\,\si{\meter}.
\]

Determine:

\begin{enumerate}[label=\alph*)]
    \item El vector posición de \(A\) respecto de \(B\).
    \item La distancia entre ambas partículas.
    \item La dirección de \(A\) respecto de \(B\).
\end{enumerate}
\end{ejercicio}

\begin{respuesta}
\[
\boxed{
\vect{r}_{A/B}
=
(8.0\ihat-6.0\jhat)\,\si{\meter}
}
\]

\[
\boxed{
r_{A/B}
=
\SI{10.0}{\meter}
}
\]

La dirección es \(36.9\grado\) por debajo de \(+x\).
\end{respuesta}

\begin{explicacion}
Elegimos el origen y los ejes cartesianos
del sistema de referencia dado.

La posición relativa es:

\[
\vect{r}_{A/B}
=
\vect{r}_A-\vect{r}_B.
\]

Por componentes:

\[
x_{A/B}=12.0-4.0=\SI{8.0}{\meter},
\]

\[
y_{A/B}=8.0-14.0=-\SI{6.0}{\meter}.
\]

Por tanto:

\[
\boxed{
\vect{r}_{A/B}
=
(8.0\ihat-6.0\jhat)\,\si{\meter}
}.
\]

La distancia entre ambas partículas es
la magnitud de este vector:

\[
r_{A/B}
=
\sqrt{(8.0)^2+(-6.0)^2}\,\si{\meter}.
\]

Así:

\[
\boxed{
r_{A/B}=\SI{10.0}{\meter}
}.
\]

Para la dirección observamos:

\[
x_{A/B}>0,
\qquad
y_{A/B}<0.
\]

Por tanto, el vector pertenece al cuarto cuadrante.

El ángulo de referencia es:

\[
\phi
=
\arctan\left(\frac{6.0}{8.0}\right).
\]

Entonces:

\[
\boxed{
\phi\approx36.9\grado
}.
\]

La partícula \(A\) se encuentra
\(36.9\grado\) por debajo del eje \(+x\),
vista desde \(B\).
\end{explicacion}


\begin{ejercicio}
Dos partículas tienen aceleraciones:

\[
\vect{a}_A
=
(3.0\ihat-2.0\jhat)\,
\si{\meter\per\second\squared},
\]

\[
\vect{a}_B
=
(-1.0\ihat+4.0\jhat)\,
\si{\meter\per\second\squared}.
\]

Determine:

\begin{enumerate}[label=\alph*)]
    \item La aceleración de \(A\) respecto de \(B\).
    \item La aceleración de \(B\) respecto de \(A\).
    \item La magnitud de la aceleración relativa.
\end{enumerate}
\end{ejercicio}

\begin{respuesta}
\[
\boxed{
\vect{a}_{A/B}
=
(4.0\ihat-6.0\jhat)\,
\si{\meter\per\second\squared}
}
\]

\[
\boxed{
\vect{a}_{B/A}
=
(-4.0\ihat+6.0\jhat)\,
\si{\meter\per\second\squared}
}
\]

\[
\boxed{
a_{A/B}\approx
\SI{7.21}{\meter\per\second\squared}
}
\]
\end{respuesta}

\begin{explicacion}
La aceleración relativa se obtiene restando
los vectores de aceleración medidos en el mismo
sistema de referencia:

\[
\vect{a}_{A/B}
=
\vect{a}_A-\vect{a}_B.
\]

Por componentes:

\[
a_{A/B,x}=3.0-(-1.0)=4.0,
\]

\[
a_{A/B,y}=-2.0-4.0=-6.0.
\]

Así:

\[
\boxed{
\vect{a}_{A/B}
=
(4.0\ihat-6.0\jhat)\,
\si{\meter\per\second\squared}
}.
\]

Al invertir el orden de las partículas:

\[
\vect{a}_{B/A}
=
-\vect{a}_{A/B}.
\]

Por tanto:

\[
\boxed{
\vect{a}_{B/A}
=
(-4.0\ihat+6.0\jhat)\,
\si{\meter\per\second\squared}
}.
\]

La magnitud es:

\[
a_{A/B}
=
\sqrt{4.0^2+(-6.0)^2}\,
\si{\meter\per\second\squared}.
\]

Por tanto:

\[
\boxed{
a_{A/B}\approx
\SI{7.21}{\meter\per\second\squared}
}.
\]

Esta aceleración representa la tasa de cambio
de la velocidad relativa, no la aceleración
absoluta de ninguna de las partículas.
\end{explicacion}


\begin{ejercicio}
Considere las siguientes afirmaciones:

\begin{enumerate}[label=\alph*)]
    \item Un objeto puede estar en reposo respecto
    de un observador y moverse respecto de otro.
    \item Si dos partículas tienen la misma velocidad
    vectorial, su velocidad relativa es cero.
    \item Si dos partículas tienen aceleraciones iguales,
    sus velocidades necesariamente son iguales.
    \item Las velocidades relativas
    \(\vect{v}_{A/B}\) y \(\vect{v}_{B/A}\)
    tienen la misma magnitud.
    \item La rapidez relativa siempre representa
    la rapidez de acercamiento entre dos partículas.
\end{enumerate}

Determine cuáles son verdaderas y justifique.
\end{ejercicio}

\begin{respuesta}
Son verdaderas:

\[
\boxed{\text{a), b) y d)}}.
\]

Son falsas:

\[
\boxed{\text{c) y e)}}.
\]
\end{respuesta}

\begin{explicacion}
\textbf{a) Verdadera.}

Una persona sentada dentro de un tren está
en reposo respecto del tren, pero se mueve
respecto del suelo cuando el tren avanza.

\textbf{b) Verdadera.}

Si:

\[
\vect{v}_A=\vect{v}_B,
\]

entonces:

\[
\vect{v}_{A/B}
=
\vect{v}_A-\vect{v}_B
=
\vect{0}.
\]

\textbf{c) Falsa.}

Si:

\[
\vect{a}_A=\vect{a}_B,
\]

se cumple:

\[
\vect{a}_{A/B}=\vect{0}.
\]

Esto implica velocidad relativa constante,
pero no necesariamente nula.

\textbf{d) Verdadera.}

Como:

\[
\vect{v}_{A/B}
=
-\vect{v}_{B/A},
\]

sus magnitudes son iguales.

\textbf{e) Falsa.}

En varias dimensiones, la velocidad relativa
puede tener una componente perpendicular
a la separación entre partículas.

Esa componente cambia la dirección relativa,
pero no contribuye instantáneamente al cambio
de distancia.

Por ello, rapidez relativa y rapidez de
acercamiento no son necesariamente iguales.
\end{explicacion}


\nivelintermedio

\begin{ejercicio}
Una banda transportadora se mueve hacia el este
con rapidez \(\SI{2.50}{\meter\per\second}\).

Una persona camina sobre ella hacia el norte
con rapidez \(\SI{1.80}{\meter\per\second}\),
medida respecto de la banda.

Determine el vector velocidad de la persona
respecto del suelo, su rapidez y su dirección.
\end{ejercicio}

\begin{respuesta}
\[
\boxed{
\vect{v}_{P/S}
=
(2.50\ihat+1.80\jhat)\,
\si{\meter\per\second}
}
\]

\[
\boxed{
v_{P/S}\approx
\SI{3.08}{\meter\per\second}
}
\]

Dirección:

\[
\boxed{
35.8\grado
}
\]

al norte del este.
\end{respuesta}

\begin{explicacion}
Elegimos:

\[
+x:\text{ este},
\qquad
+y:\text{ norte}.
\]

La velocidad de la banda respecto del suelo es:

\[
\vect{v}_{B/S}
=
2.50\ihat\,
\si{\meter\per\second}.
\]

La velocidad de la persona respecto
de la banda es:

\[
\vect{v}_{P/B}
=
1.80\jhat\,
\si{\meter\per\second}.
\]

La composición vectorial es:

\[
\vect{v}_{P/S}
=
\vect{v}_{P/B}
+
\vect{v}_{B/S}.
\]

Por componentes:

\[
v_{P/S,x}=2.50,
\qquad
v_{P/S,y}=1.80.
\]

Por tanto:

\[
\boxed{
\vect{v}_{P/S}
=
(2.50\ihat+1.80\jhat)\,
\si{\meter\per\second}
}.
\]

La magnitud es:

\[
v_{P/S}
=
\sqrt{2.50^2+1.80^2}.
\]

Entonces:

\[
\boxed{
v_{P/S}\approx
\SI{3.08}{\meter\per\second}
}.
\]

Como ambas componentes son positivas,
la dirección está en el primer cuadrante.

El ángulo respecto del este es:

\[
\theta
=
\arctan\left(\frac{1.80}{2.50}\right).
\]

Así:

\[
\boxed{
\theta\approx35.8\grado
}.
\]

Para un observador fijo al suelo, la persona
se mueve hacia el noreste, aunque respecto
de la banda camina directamente hacia el norte.
\end{explicacion}


\begin{ejercicio}
Dos automóviles recorren una carretera recta
con velocidades constantes.

En el instante inicial, \(A\) se encuentra
en \(x_{A0}=\SI{0}{\meter}\) y \(B\)
en \(x_{B0}=\SI{180}{\meter}\).

Ambos avanzan hacia \(+x\), con:

\[
v_A=\SI{28.0}{\meter\per\second},
\qquad
v_B=\SI{18.0}{\meter\per\second}.
\]

Determine el instante y la posición
en que \(A\) alcanza a \(B\).

Resuelva el problema mediante posición relativa
y compruebe el resultado con las ecuaciones
individuales del movimiento.
\end{ejercicio}

\begin{respuesta}
\[
\boxed{
t=\SI{18.0}{\second}
}
\]

\[
\boxed{
x=\SI{504}{\meter}
}
\]
\end{respuesta}

\begin{explicacion}
Tomamos el origen en la posición inicial
de \(A\), con \(+x\) en el sentido de avance.

La posición relativa inicial es:

\[
x_{A/B,0}
=
x_{A0}-x_{B0}.
\]

Por tanto:

\[
x_{A/B,0}=0-180=-\SI{180}{\meter}.
\]

La velocidad relativa es:

\[
v_{A/B}=v_A-v_B.
\]

Entonces:

\[
v_{A/B}=28.0-18.0
=\SI{10.0}{\meter\per\second}.
\]

Como ambos movimientos son uniformes:

\[
x_{A/B}(t)
=
x_{A/B,0}+v_{A/B}t.
\]

Los vehículos se encuentran cuando:

\[
x_{A/B}(t)=0.
\]

Así:

\[
0=-180+10.0t.
\]

Despejando:

\[
\boxed{
t=\SI{18.0}{\second}
}.
\]

La posición de \(A\) es:

\[
x_A(t)=x_{A0}+v_At.
\]

Entonces:

\[
x_A=0+28.0(18.0).
\]

Por tanto:

\[
\boxed{
x_A=\SI{504}{\meter}
}.
\]

Comprobamos con \(B\):

\[
x_B=180+18.0(18.0).
\]

Así:

\[
x_B=\SI{504}{\meter}.
\]

Como:

\[
x_A=x_B,
\]

el encuentro queda verificado.

El movimiento relativo permite convertir
dos movimientos en uno solo: el cierre
de la separación inicial.
\end{explicacion}


\begin{ejercicio}
Una embarcación cruza un río de ancho
\(D=\SI{120}{\meter}\).

Su velocidad respecto del agua es perpendicular
a las orillas, con rapidez
\(\SI{4.00}{\meter\per\second}\).

La corriente uniforme se desplaza aguas abajo
con rapidez \(\SI{2.50}{\meter\per\second}\).

Determine:

\begin{enumerate}[label=\alph*)]
    \item La velocidad de la embarcación respecto del suelo.
    \item El tiempo de cruce.
    \item El desplazamiento aguas abajo.
    \item La rapidez respecto del suelo.
    \item El ángulo de la trayectoria respecto de
    la perpendicular a las orillas.
\end{enumerate}
\end{ejercicio}

\begin{respuesta}
\[
\boxed{
\vect{v}_{B/S}
=
(4.00\ihat+2.50\jhat)\,
\si{\meter\per\second}
}
\]

\[
\boxed{
t=\SI{30.0}{\second}
}
\]

\[
\boxed{
\Delta y=\SI{75.0}{\meter}
}
\]

\[
\boxed{
v_{B/S}\approx
\SI{4.72}{\meter\per\second}
}
\]

\[
\boxed{
\theta\approx32.0\grado
}
\]

aguas abajo respecto de la perpendicular.
\end{respuesta}

\begin{explicacion}
Elegimos:

\[
+x:\text{ hacia la orilla opuesta},
\]

\[
+y:\text{ aguas abajo}.
\]

La velocidad respecto del agua es:

\[
\vect{v}_{B/A}
=
4.00\ihat\,
\si{\meter\per\second}.
\]

La corriente tiene:

\[
\vect{v}_{A/S}
=
2.50\jhat\,
\si{\meter\per\second}.
\]

La velocidad respecto del suelo es:

\[
\vect{v}_{B/S}
=
\vect{v}_{B/A}+\vect{v}_{A/S}.
\]

Por tanto:

\[
\boxed{
\vect{v}_{B/S}
=
(4.00\ihat+2.50\jhat)\,
\si{\meter\per\second}
}.
\]

La componente perpendicular determina el
tiempo de cruce:

\[
D=v_{B/S,x}t.
\]

Despejando:

\[
t=\frac{D}{v_{B/S,x}}.
\]

Sustituyendo:

\[
t
=
\frac{\SI{120}{\meter}}
{\SI{4.00}{\meter\per\second}}.
\]

Por tanto:

\[
\boxed{
t=\SI{30.0}{\second}
}.
\]

El desplazamiento aguas abajo es:

\[
\Delta y=v_{B/S,y}t.
\]

Entonces:

\[
\Delta y=(2.50)(30.0).
\]

Así:

\[
\boxed{
\Delta y=\SI{75.0}{\meter}
}.
\]

La rapidez respecto del suelo es:

\[
v_{B/S}
=
\sqrt{4.00^2+2.50^2}.
\]

Por tanto:

\[
\boxed{
v_{B/S}\approx
\SI{4.72}{\meter\per\second}
}.
\]

El ángulo de la trayectoria satisface:

\[
\tan\theta
=
\frac{v_{B/S,y}}{v_{B/S,x}}.
\]

Entonces:

\[
\theta
=
\arctan\left(\frac{2.50}{4.00}\right).
\]

Así:

\[
\boxed{
\theta\approx32.0\grado
}.
\]

La embarcación apunta directamente hacia
la otra orilla, pero la corriente hace
que su trayectoria respecto del suelo
se incline aguas abajo.
\end{explicacion}


\begin{ejercicio}
Una embarcación debe cruzar un río de ancho
\(D=\SI{160}{\meter}\) y llegar exactamente
al punto situado frente al lugar de partida.

La corriente uniforme tiene rapidez
\(\SI{3.00}{\meter\per\second}\).

La embarcación mantiene una rapidez constante
de \(\SI{5.00}{\meter\per\second}\)
respecto del agua.

Determine:

\begin{enumerate}[label=\alph*)]
    \item Las componentes de la velocidad
    de la embarcación respecto del agua.
    \item El ángulo que debe formar su rumbo
    con la perpendicular a las orillas.
    \item Su velocidad respecto del suelo.
    \item El tiempo de cruce.
\end{enumerate}
\end{ejercicio}

\begin{respuesta}
\[
\boxed{
\vect{v}_{B/A}
=
(4.00\ihat-3.00\jhat)\,
\si{\meter\per\second}
}
\]

\[
\boxed{
\phi\approx36.9\grado
}
\]

hacia aguas arriba desde la perpendicular.

\[
\boxed{
\vect{v}_{B/S}
=
4.00\ihat\,
\si{\meter\per\second}
}
\]

\[
\boxed{
t=\SI{40.0}{\second}
}
\]
\end{respuesta}

\begin{explicacion}
Elegimos \(+x\) hacia la orilla opuesta
y \(+y\) aguas abajo.

La corriente tiene velocidad:

\[
\vect{v}_{A/S}
=
3.00\jhat\,
\si{\meter\per\second}.
\]

Para llegar directamente al punto opuesto,
la velocidad respecto del suelo no puede
tener componente aguas abajo:

\[
v_{B/S,y}=0.
\]

Por composición de velocidades:

\[
v_{B/S,y}
=
v_{B/A,y}+v_{A/S,y}.
\]

Entonces:

\[
0=v_{B/A,y}+3.00.
\]

Por tanto:

\[
\boxed{
v_{B/A,y}=-\SI{3.00}{\meter\per\second}
}.
\]

La rapidez respecto del agua es:

\[
v_{B/A}=\SI{5.00}{\meter\per\second}.
\]

Por tanto:

\[
v_{B/A}^2
=
v_{B/A,x}^2+v_{B/A,y}^2.
\]

Despejando la componente perpendicular:

\[
v_{B/A,x}
=
\sqrt{v_{B/A}^2-v_{B/A,y}^2}.
\]

Entonces:

\[
v_{B/A,x}
=
\sqrt{5.00^2-3.00^2}.
\]

Así:

\[
\boxed{
v_{B/A,x}=\SI{4.00}{\meter\per\second}
}.
\]

El vector velocidad respecto del agua es:

\[
\boxed{
\vect{v}_{B/A}
=
(4.00\ihat-3.00\jhat)\,
\si{\meter\per\second}
}.
\]

El ángulo respecto de la perpendicular
a las orillas es:

\[
\tan\phi=\frac{3.00}{4.00}.
\]

Por tanto:

\[
\boxed{
\phi\approx36.9\grado
}
\]

hacia aguas arriba.

La velocidad respecto del suelo es:

\[
\vect{v}_{B/S}
=
(4.00\ihat-3.00\jhat)
+
3.00\jhat.
\]

Entonces:

\[
\boxed{
\vect{v}_{B/S}
=
4.00\ihat\,
\si{\meter\per\second}
}.
\]

Finalmente:

\[
t=\frac{D}{v_{B/S,x}}
=
\frac{160}{4.00}.
\]

Por tanto:

\[
\boxed{
t=\SI{40.0}{\second}
}.
\]

El ángulo hacia aguas arriba cancela la
deriva producida por la corriente.
\end{explicacion}


\begin{ejercicio}
Una aeronave debe desplazarse directamente
hacia el norte.

Su rapidez respecto del aire es
\(\SI{90.0}{\meter\per\second}\).

El viento sopla hacia el este con rapidez
\(\SI{30.0}{\meter\per\second}\).

Determine:

\begin{enumerate}[label=\alph*)]
    \item Las componentes necesarias de la velocidad
    de la aeronave respecto del aire.
    \item El ángulo de corrección del rumbo.
    \item La rapidez respecto del suelo.
    \item El tiempo necesario para recorrer
    \(\SI{45.0}{\kilo\meter}\) hacia el norte,
    suponiendo condiciones constantes.
\end{enumerate}
\end{ejercicio}

\begin{respuesta}
\[
\boxed{
\vect{v}_{P/A}
\approx
(-30.0\ihat+84.9\jhat)\,
\si{\meter\per\second}
}
\]

\[
\boxed{
\phi\approx19.5\grado
}
\]

hacia el oeste desde el norte.

\[
\boxed{
v_{P/S}\approx
\SI{84.9}{\meter\per\second}
}
\]

\[
\boxed{
t\approx\SI{530}{\second}
\approx\SI{8.84}{\minute}
}
\]
\end{respuesta}

\begin{explicacion}
Elegimos:

\[
+x:\text{ este},
\qquad
+y:\text{ norte}.
\]

La velocidad del viento es:

\[
\vect{v}_{A/S}
=
30.0\ihat\,
\si{\meter\per\second}.
\]

Para avanzar directamente hacia el norte,
la velocidad respecto del suelo debe cumplir:

\[
v_{P/S,x}=0.
\]

Aplicamos:

\[
v_{P/S,x}
=
v_{P/A,x}+v_{A/S,x}.
\]

Entonces:

\[
0=v_{P/A,x}+30.0.
\]

Por tanto:

\[
\boxed{
v_{P/A,x}
=
-\SI{30.0}{\meter\per\second}
}.
\]

La rapidez respecto del aire satisface:

\[
v_{P/A}^2
=
v_{P/A,x}^2+v_{P/A,y}^2.
\]

Despejando:

\[
v_{P/A,y}
=
\sqrt{90.0^2-30.0^2}.
\]

Entonces:

\[
\boxed{
v_{P/A,y}
\approx
\SI{84.9}{\meter\per\second}
}.
\]

Por tanto:

\[
\boxed{
\vect{v}_{P/A}
\approx
(-30.0\ihat+84.9\jhat)\,
\si{\meter\per\second}
}.
\]

El ángulo hacia el oeste desde el norte
satisface:

\[
\sen\phi=\frac{30.0}{90.0}.
\]

Entonces:

\[
\boxed{
\phi\approx19.5\grado
}.
\]

Al sumar el viento:

\[
\vect{v}_{P/S}
=
(-30.0\ihat+84.9\jhat)
+
30.0\ihat.
\]

Así:

\[
\boxed{
\vect{v}_{P/S}
\approx
84.9\jhat\,
\si{\meter\per\second}
}.
\]

La distancia hacia el norte es:

\[
D=\SI{45.0}{\kilo\meter}
=
\SI{45000}{\meter}.
\]

Como la velocidad resultante es constante:

\[
t=\frac{D}{v_{P/S,y}}.
\]

Por tanto:

\[
t
=
\frac{45000}{\sqrt{7200}}.
\]

Así:

\[
\boxed{
t\approx\SI{530}{\second}
\approx\SI{8.84}{\minute}
}.
\]

La aeronave apunta hacia el noroeste
para compensar el viento del este,
pero su trayectoria respecto del suelo
es directamente hacia el norte.
\end{explicacion}


\begin{ejercicio}
Dos partículas se desplazan con velocidades
constantes en un plano.

En \(t=0\), sus posiciones son:

\[
\vect{r}_{A0}
=
(0\ihat+30.0\jhat)\,\si{\meter},
\]

\[
\vect{r}_{B0}
=
(40.0\ihat+0\jhat)\,\si{\meter}.
\]

Sus velocidades son:

\[
\vect{v}_A
=
(6.0\ihat+1.0\jhat)\,
\si{\meter\per\second},
\]

\[
\vect{v}_B
=
(2.0\ihat+4.0\jhat)\,
\si{\meter\per\second}.
\]

Determine si las partículas se encuentran
y, en caso afirmativo, el instante
y la posición del encuentro.
\end{ejercicio}

\begin{respuesta}
Las partículas se encuentran en:

\[
\boxed{
t=\SI{10.0}{\second}
}
\]

en la posición:

\[
\boxed{
\vect{r}
=
(60.0\ihat+40.0\jhat)\,\si{\meter}
}.
\]
\end{respuesta}

\begin{explicacion}
El modelo corresponde a dos movimientos
rectilíneos uniformes en el plano.

La posición relativa inicial es:

\[
\vect{r}_{A/B,0}
=
\vect{r}_{A0}-\vect{r}_{B0}.
\]

Entonces:

\[
\vect{r}_{A/B,0}
=
(-40.0\ihat+30.0\jhat)\,\si{\meter}.
\]

La velocidad relativa es:

\[
\vect{v}_{A/B}
=
\vect{v}_A-\vect{v}_B.
\]

Por componentes:

\[
v_{A/B,x}=6.0-2.0=4.0,
\]

\[
v_{A/B,y}=1.0-4.0=-3.0.
\]

Por tanto:

\[
\vect{v}_{A/B}
=
(4.0\ihat-3.0\jhat)\,
\si{\meter\per\second}.
\]

La posición relativa es:

\[
\vect{r}_{A/B}(t)
=
(-40.0+4.0t)\ihat
+
(30.0-3.0t)\jhat.
\]

Para que exista encuentro:

\[
\vect{r}_{A/B}(t)=\vect{0}.
\]

Esto exige:

\[
-40.0+4.0t=0,
\]

\[
30.0-3.0t=0.
\]

Ambas ecuaciones proporcionan:

\[
\boxed{
t=\SI{10.0}{\second}
}.
\]

Determinamos la posición utilizando
el movimiento de \(A\):

\[
\vect{r}_A(t)
=
\vect{r}_{A0}+\vect{v}_At.
\]

Entonces:

\[
\vect{r}_A(10.0)
=
(0\ihat+30.0\jhat)
+
(6.0\ihat+1.0\jhat)(10.0).
\]

Por tanto:

\[
\boxed{
\vect{r}
=
(60.0\ihat+40.0\jhat)\,\si{\meter}
}.
\]

Comprobando con \(B\):

\[
\vect{r}_B(10.0)
=
(40.0\ihat+0\jhat)
+
(2.0\ihat+4.0\jhat)(10.0).
\]

Se obtiene la misma posición.

Como coinciden ambas coordenadas
en el mismo instante, existe encuentro.
\end{explicacion}


\nivelavanzado

\begin{ejercicio}
Dos vehículos se desplazan en un plano
con velocidades constantes.

En el instante \(t=0\), la posición
de \(A\) respecto de \(B\) es:

\[
\vect{r}_{A/B,0}
=
(100\ihat+60\jhat)\,\si{\meter}.
\]

La velocidad relativa es:

\[
\vect{v}_{A/B}
=
(-10\ihat-2\jhat)\,
\si{\meter\per\second}.
\]

Determine:

\begin{enumerate}[label=\alph*)]
    \item Si inicialmente se acercan o se alejan.
    \item El instante en que alcanzan
    su separación mínima.
    \item La separación mínima.
    \item Si llegan a encontrarse.
\end{enumerate}
\end{ejercicio}

\begin{respuesta}
Inicialmente se acercan.

\[
\boxed{
t_{\min}\approx
\SI{10.77}{\second}
}
\]

\[
\boxed{
d_{\min}\approx
\SI{39.2}{\meter}
}
\]

No llegan a encontrarse.
\end{respuesta}

\begin{explicacion}
La posición relativa para velocidad
constante es:

\[
\vect{r}(t)=\vect{r}_0+\vect{v}t,
\]

donde:

\[
\vect{r}_0
=
(100\ihat+60\jhat)\,\si{\meter},
\]

\[
\vect{v}
=
(-10\ihat-2\jhat)\,
\si{\meter\per\second}.
\]

Para saber si inicialmente se acercan,
analizamos:

\[
\vect{r}_0\pdot\vect{v}.
\]

Entonces:

\[
\vect{r}_0\pdot\vect{v}
=
100(-10)+60(-2).
\]

Por tanto:

\[
\vect{r}_0\pdot\vect{v}
=
-1120\,
\si{\meter\squared\per\second}.
\]

El resultado es negativo.

Como:

\[
\frac{\mathrm{d}d}{\mathrm{d}t}
=
\frac{\vect{r}\pdot\vect{v}}{\norm{\vect{r}}},
\]

la distancia inicial está disminuyendo.

Para determinar la mínima separación,
consideramos:

\[
d^2(t)
=
(\vect{r}_0+\vect{v}t)
\pdot
(\vect{r}_0+\vect{v}t).
\]

Derivamos respecto del tiempo y
buscamos el mínimo:

\[
\frac{\mathrm{d}(d^2)}{\mathrm{d}t}
=
2\vect{r}_0\pdot\vect{v}
+
2v^2t=0.
\]

Entonces:

\[
t_{\min}
=
-\frac{\vect{r}_0\pdot\vect{v}}{v^2}.
\]

La rapidez relativa al cuadrado es:

\[
v^2=(-10)^2+(-2)^2=104.
\]

Por tanto:

\[
t_{\min}
=
-\frac{-1120}{104}.
\]

Así:

\[
\boxed{
t_{\min}\approx
\SI{10.77}{\second}
}.
\]

En ese instante:

\[
x_{A/B}
=
100-10(10.76923)
\approx-7.692,
\]

\[
y_{A/B}
=
60-2(10.76923)
\approx38.462.
\]

Por tanto:

\[
d_{\min}
=
\sqrt{(-7.692)^2+(38.462)^2}.
\]

Entonces:

\[
\boxed{
d_{\min}\approx
\SI{39.2}{\meter}
}.
\]

Como la distancia mínima es distinta
de cero, no existe encuentro.

El ejercicio demuestra que acercarse
no implica necesariamente encontrarse.
\end{explicacion}


\begin{ejercicio}
Una embarcación puede mantener una rapidez
\(v\) respecto del agua al cruzar un río
de ancho \(D\).

La corriente es uniforme, paralela a las
orillas y tiene rapidez \(u\), con:

\[
0<u<v.
\]

Considere dos estrategias:

\begin{enumerate}[label=\alph*)]
    \item Orientar la embarcación perpendicularmente
    a las orillas para minimizar el tiempo de cruce.
    \item Orientarla aguas arriba de manera
    que llegue exactamente al punto opuesto.
\end{enumerate}

Deduzca los tiempos de cruce de ambas estrategias
y demuestre cuál es menor.

Analice además el comportamiento cuando
\(u\) se aproxima a \(v\).
\end{ejercicio}

\begin{respuesta}
Para el cruce de menor tiempo:

\[
\boxed{
t_{\min}=\frac{D}{v}
}
\]

Para el cruce sin deriva:

\[
\boxed{
t_{\mathrm{sin\ deriva}}
=
\frac{D}{\sqrt{v^2-u^2}}
}
\]

Como \(0<u<v\):

\[
\boxed{
t_{\mathrm{sin\ deriva}}>t_{\min}
}.
\]

Cuando \(u\to v^{-}\), el tiempo
de cruce sin deriva tiende a infinito.
\end{respuesta}

\begin{explicacion}
Elegimos:

\[
+x:\text{ perpendicular hacia la orilla opuesta},
\]

\[
+y:\text{ aguas abajo}.
\]

La corriente tiene velocidad:

\[
\vect{v}_{A/S}=u\jhat.
\]

La velocidad de la embarcación respecto
del agua puede escribirse como:

\[
\vect{v}_{B/A}
=
v\cos\phi\,\ihat
-
v\sen\phi\,\jhat,
\]

donde \(\phi\) es el ángulo de orientación
hacia aguas arriba desde la perpendicular.

La velocidad respecto del suelo es:

\[
\vect{v}_{B/S}
=
v\cos\phi\,\ihat
+
(u-v\sen\phi)\jhat.
\]

El tiempo de cruce depende de la componente
perpendicular:

\[
t=\frac{D}{v\cos\phi}.
\]

\textbf{Primera estrategia: minimizar el tiempo.}

Para minimizar \(t\), debemos maximizar:

\[
v\cos\phi.
\]

Como \(v\) es constante, su valor máximo es \(v\),
obtenido cuando:

\[
\phi=0.
\]

Por tanto:

\[
\boxed{
t_{\min}=\frac{D}{v}
}.
\]

\textbf{Segunda estrategia: evitar la deriva.}

Para llegar al punto opuesto:

\[
v_{B/S,y}=0.
\]

Entonces:

\[
u-v\sen\phi=0.
\]

Por tanto:

\[
\sen\phi=\frac{u}{v}.
\]

La componente perpendicular será:

\[
v_{B/S,x}
=
v\cos\phi
=
\sqrt{v^2-u^2}.
\]

El tiempo es:

\[
\boxed{
t_{\mathrm{sin\ deriva}}
=
\frac{D}{\sqrt{v^2-u^2}}
}.
\]

Como:

\[
0<u<v,
\]

tenemos:

\[
v^2-u^2<v^2.
\]

Por tanto:

\[
\sqrt{v^2-u^2}<v.
\]

Al dividir la misma distancia \(D>0\)
entre una rapidez menor, se obtiene
un tiempo mayor:

\[
\boxed{
t_{\mathrm{sin\ deriva}}>t_{\min}
}.
\]

Finalmente, cuando:

\[
u\to v^{-},
\]

se cumple:

\[
\sqrt{v^2-u^2}\to0^{+}.
\]

En consecuencia:

\[
\boxed{
t_{\mathrm{sin\ deriva}}\to\infty
}.
\]

Físicamente, la embarcación debe orientar
casi toda su capacidad de avance aguas arriba
para compensar la corriente, quedando una
componente transversal cada vez menor.

Las dos estrategias optimizan objetivos
diferentes y no deben confundirse.
\end{explicacion}


\begin{ejercicio}
Un avión dispone de una rapidez constante
\(v\) respecto del aire.

Debe recorrer una distancia \(D\)
directamente hacia el norte.

Durante el trayecto existe viento
uniforme hacia el este, de rapidez \(u\),
con:

\[
0<u<v.
\]

Deduzca:

\begin{enumerate}[label=\alph*)]
    \item El ángulo de corrección del rumbo.
    \item La rapidez respecto del suelo
    en dirección norte.
    \item El tiempo de vuelo.
    \item El comportamiento de los resultados
    cuando \(u=0\) y cuando \(u\to v^{-}\).
\end{enumerate}
\end{ejercicio}

\begin{respuesta}
\[
\boxed{
\phi=\arcsen\left(\frac{u}{v}\right)
}
\]

\[
\boxed{
v_{P/S}=\sqrt{v^2-u^2}
}
\]

\[
\boxed{
t=\frac{D}{\sqrt{v^2-u^2}}
}
\]

Cuando \(u=0\), el avión apunta al norte
y \(t=D/v\).

Cuando \(u\to v^{-}\), la rapidez norte
tiende a cero y el tiempo tiende a infinito.
\end{respuesta}

\begin{explicacion}
Elegimos \(+x\) hacia el este
y \(+y\) hacia el norte.

El viento tiene velocidad:

\[
\vect{v}_{A/S}=u\ihat.
\]

El avión debe orientar su velocidad
respecto del aire parcialmente hacia el oeste.

Escribimos:

\[
\vect{v}_{P/A}
=
-v\sen\phi\,\ihat
+
v\cos\phi\,\jhat,
\]

donde \(\phi\) se mide hacia el oeste
desde el norte.

La velocidad respecto del suelo es:

\[
\vect{v}_{P/S}
=
\vect{v}_{P/A}
+
\vect{v}_{A/S}.
\]

Por componentes:

\[
v_{P/S,x}
=
-v\sen\phi+u,
\]

\[
v_{P/S,y}
=
v\cos\phi.
\]

Para seguir una trayectoria directamente
hacia el norte:

\[
v_{P/S,x}=0.
\]

Por tanto:

\[
-v\sen\phi+u=0.
\]

Despejando:

\[
\sen\phi=\frac{u}{v}.
\]

Así:

\[
\boxed{
\phi=\arcsen\left(\frac{u}{v}\right)
}.
\]

La rapidez norte es:

\[
v_{P/S}
=
v\cos\phi.
\]

Como:

\[
\cos\phi
=
\sqrt{1-\sen^2\phi},
\]

obtenemos:

\[
v_{P/S}
=
v\sqrt{1-\frac{u^2}{v^2}}.
\]

Simplificando:

\[
\boxed{
v_{P/S}=\sqrt{v^2-u^2}
}.
\]

Para recorrer la distancia \(D\):

\[
D=v_{P/S}t.
\]

Por tanto:

\[
\boxed{
t=\frac{D}{\sqrt{v^2-u^2}}
}.
\]

Si desaparece el viento:

\[
u=0,
\]

entonces:

\[
\phi=0,
\qquad
v_{P/S}=v,
\qquad
t=\frac{D}{v}.
\]

En cambio, cuando:

\[
u\to v^{-},
\]

se obtiene:

\[
\phi\to90\grado,
\]

\[
v_{P/S}\to0,
\]

\[
t\to\infty.
\]

Esto significa que casi toda la velocidad
del avión respecto del aire debe utilizarse
para compensar el viento lateral.

El modelo exige \(u<v\) para mantener
una componente norte positiva.
\end{explicacion}


\begin{ejercicio}
Una partícula \(P\) se mueve en un sistema
de referencia \(S\) con posición:

\[
\vect{r}_P(t)
=
(4t+2)\ihat+(3t^2-1)\jhat,
\]

donde las coordenadas se expresan en metros
y el tiempo en segundos.

Un sistema móvil \(S'\) tiene su origen
en la posición:

\[
\vect{R}(t)
=
(2t-3)\ihat+(t+4)\jhat.
\]

Los ejes de ambos sistemas permanecen paralelos.

Determine:

\begin{enumerate}[label=\alph*)]
    \item La posición de \(P\) respecto de \(S'\).
    \item La velocidad respecto de ambos sistemas.
    \item La aceleración respecto de ambos sistemas.
    \item La posición y velocidad relativas
    cuando \(t=\SI{2.00}{\second}\).
\end{enumerate}
\end{ejercicio}

\begin{respuesta}
\[
\boxed{
\vect{r}'_P(t)
=
(2t+5)\ihat+(3t^2-t-5)\jhat
}
\]

\[
\boxed{
\vect{v}_{P/S}
=
4\ihat+6t\jhat
}
\]

\[
\boxed{
\vect{v}_{P/S'}
=
2\ihat+(6t-1)\jhat
}
\]

\[
\boxed{
\vect{a}_{P/S}
=
\vect{a}_{P/S'}
=
6\jhat
}
\]

Con \(t=\SI{2.00}{\second}\):

\[
\boxed{
\vect{r}'_P
=
(9\ihat+5\jhat)\,\si{\meter}
}
\]

\[
\boxed{
\vect{v}_{P/S'}
=
(2\ihat+11\jhat)\,
\si{\meter\per\second}
}
\]
\end{respuesta}

\begin{explicacion}
La posición en el sistema fijo satisface:

\[
\vect{r}_P
=
\vect{R}+\vect{r}'_P.
\]

Por tanto:

\[
\vect{r}'_P
=
\vect{r}_P-\vect{R}.
\]

Trabajando por componentes:

\[
x'=4t+2-(2t-3),
\]

\[
y'=3t^2-1-(t+4).
\]

Simplificando:

\[
\boxed{
\vect{r}'_P(t)
=
(2t+5)\ihat+(3t^2-t-5)\jhat
}.
\]

Ahora derivamos las posiciones
respecto del tiempo.

En el sistema \(S\):

\[
\vect{v}_{P/S}
=
\frac{\mathrm{d}\vect{r}_P}{\mathrm{d}t}.
\]

Entonces:

\[
\boxed{
\vect{v}_{P/S}
=
4\ihat+6t\jhat
}.
\]

La velocidad del origen móvil es:

\[
\vect{v}_{S'/S}
=
\frac{\mathrm{d}\vect{R}}{\mathrm{d}t}.
\]

Por tanto:

\[
\vect{v}_{S'/S}
=
2\ihat+\jhat.
\]

La velocidad relativa es:

\[
\vect{v}_{P/S'}
=
\vect{v}_{P/S}-\vect{v}_{S'/S}.
\]

Entonces:

\[
\boxed{
\vect{v}_{P/S'}
=
2\ihat+(6t-1)\jhat
}.
\]

Derivando las velocidades:

\[
\boxed{
\vect{a}_{P/S}=6\jhat\,
\si{\meter\per\second\squared}
},
\]

\[
\boxed{
\vect{a}_{P/S'}=6\jhat\,
\si{\meter\per\second\squared}
}.
\]

Las aceleraciones son iguales porque
el sistema \(S'\) se traslada con
velocidad vectorial constante.

Para \(t=\SI{2.00}{\second}\):

\[
x'=2(2)+5=9,
\]

\[
y'=3(2)^2-2-5=5.
\]

Por tanto:

\[
\boxed{
\vect{r}'_P
=
(9\ihat+5\jhat)\,\si{\meter}
}.
\]

La velocidad relativa es:

\[
\vect{v}_{P/S'}
=
2\ihat+[6(2)-1]\jhat.
\]

Así:

\[
\boxed{
\vect{v}_{P/S'}
=
(2\ihat+11\jhat)\,
\si{\meter\per\second}
}.
\]

Este ejercicio muestra que la posición
y la velocidad dependen del sistema
de referencia, mientras que la aceleración
es la misma en dos sistemas relacionados
por traslación uniforme.
\end{explicacion}


\begin{ejercicio}
Dos partículas \(A\) y \(B\) se mueven
con velocidades constantes en un plano.

Sus posiciones iniciales son:

\[
\vect{r}_{A0}
=
(0\ihat+20\jhat)\,\si{\meter},
\]

\[
\vect{r}_{B0}
=
(30\ihat+0\jhat)\,\si{\meter}.
\]

La partícula \(A\) posee velocidad:

\[
\vect{v}_A
=
(5\ihat+2\jhat)\,
\si{\meter\per\second}.
\]

Se sabe que ambas partículas se encuentran
exactamente en:

\[
t=\SI{10.0}{\second}.
\]

Determine:

\begin{enumerate}[label=\alph*)]
    \item La velocidad que debe tener \(B\).
    \item La velocidad de \(A\) respecto de \(B\).
    \item La posición del encuentro.
    \item Si la solución es única bajo
    el modelo de velocidades constantes.
\end{enumerate}
\end{ejercicio}

\begin{respuesta}
\[
\boxed{
\vect{v}_B
=
(2\ihat+4\jhat)\,
\si{\meter\per\second}
}
\]

\[
\boxed{
\vect{v}_{A/B}
=
(3\ihat-2\jhat)\,
\si{\meter\per\second}
}
\]

\[
\boxed{
\vect{r}_{\mathrm{encuentro}}
=
(50\ihat+40\jhat)\,\si{\meter}
}
\]

La solución es única bajo las hipótesis dadas.
\end{respuesta}

\begin{explicacion}
Para velocidades constantes:

\[
\vect{r}_A(t)
=
\vect{r}_{A0}+\vect{v}_At,
\]

\[
\vect{r}_B(t)
=
\vect{r}_{B0}+\vect{v}_Bt.
\]

La condición de encuentro es:

\[
\vect{r}_A(t)=\vect{r}_B(t).
\]

Como el instante está dado:

\[
t=\SI{10.0}{\second},
\]

podemos determinar primero la posición
de \(A\) en ese instante:

\[
\vect{r}_A(10)
=
(0\ihat+20\jhat)
+
(5\ihat+2\jhat)(10).
\]

Entonces:

\[
\boxed{
\vect{r}_{\mathrm{encuentro}}
=
(50\ihat+40\jhat)\,\si{\meter}
}.
\]

La partícula \(B\) debe llegar
a esa misma posición en el mismo tiempo.

Por tanto:

\[
\vect{r}_{B0}+\vect{v}_Bt
=
\vect{r}_{\mathrm{encuentro}}.
\]

Despejando:

\[
\vect{v}_B
=
\frac{
\vect{r}_{\mathrm{encuentro}}
-
\vect{r}_{B0}
}{t}.
\]

Por componentes:

\[
v_{Bx}
=
\frac{50-30}{10}
=
\SI{2.0}{\meter\per\second},
\]

\[
v_{By}
=
\frac{40-0}{10}
=
\SI{4.0}{\meter\per\second}.
\]

Así:

\[
\boxed{
\vect{v}_B
=
(2\ihat+4\jhat)\,
\si{\meter\per\second}
}.
\]

La velocidad relativa es:

\[
\vect{v}_{A/B}
=
\vect{v}_A-\vect{v}_B.
\]

Entonces:

\[
\vect{v}_{A/B}
=
(5-2)\ihat+(2-4)\jhat.
\]

Por tanto:

\[
\boxed{
\vect{v}_{A/B}
=
(3\ihat-2\jhat)\,
\si{\meter\per\second}
}.
\]

Para verificar, calculamos:

\[
\vect{r}_{A/B,0}
=
(-30\ihat+20\jhat)\,\si{\meter}.
\]

A los diez segundos:

\[
\vect{r}_{A/B}(10)
=
(-30\ihat+20\jhat)
+
(3\ihat-2\jhat)(10).
\]

Así:

\[
\vect{r}_{A/B}(10)=\vect{0}.
\]

La solución es única porque están fijados
la posición inicial, la posición final
y el tiempo, y se exige una velocidad
vectorial constante durante el intervalo.
\end{explicacion}